\documentclass{article}

\usepackage{amsmath,amssymb}
\usepackage{xurl}
\usepackage[hidelinks]{hyperref}

% Shared commands for notes in docs/paper-gaps/.
% Notation follows the LDT paper (references/ldt-paper/).

\newcommand{\Lean}{\textsc{Lean}}
\newcommand{\leanid}[1]{\textsf{#1}}
\newcommand{\ghissue}[1]{\href{https://github.com/LionSR/MIPStarRE/issues/#1}{GitHub issue \##1}}
\newcommand{\ghpr}[1]{\href{https://github.com/LionSR/MIPStarRE/pull/#1}{GitHub PR \##1}}

% --- Paper-compatible notation ---
\newcommand{\F}{\mathbb{F}}
\newcommand{\N}{\mathbb{N}}
\newcommand{\C}{\mathbb{C}}
\newcommand{\tr}{\operatorname{tr}}
\newcommand{\ot}{\otimes}
\newcommand{\E}{\mathop{\bf E\/}}
\newcommand{\bx}{\mathbf{x}}
\newcommand{\by}{\mathbf{y}}
\newcommand{\bu}{\mathbf{u}}
\newcommand{\bv}{\mathbf{v}}

% Approximation relation (paper uses \approx_\eps and \simeq_\zeta)
\newcommand{\approxeps}{\approx_{\varepsilon}}
\newcommand{\simgeq}{\simeq_{\zeta}}

% Polynomial spaces (paper uses \polyfunc, \polysub, \polymeas)
\newcommand{\polyfunc}[3]{\operatorname{Poly}(#1,#2,#3)}
\newcommand{\polysub}[3]{\operatorname{PolySub}(#1,#2,#3)}
\newcommand{\polymeas}[3]{\operatorname{PolyMeas}(#1,#2,#3)}


\title{Foundation-layer formalization design: $\mathbf{F}_q$, distributions, and zero families}
\author{}
\date{\today}

\begin{document}
\maketitle

\section{The assertion}

This note concerns the three foundation-layer design choices identified in
Issue~\#458 of the \Lean{} formalization of the low individual degree test
\cite{Ji2020LowIndividualDegree}. The choices are:
\begin{enumerate}
  \item The paper parameter $q$ is a prime power specifying the finite field
    $\F_q$;\footnote{\texttt{references/ldt-paper/preliminaries.tex},
      lines~17--19 and~89--93.}
    the \Lean{} structure \texttt{Parameters} stores $q$ as a bare
    $\mathbb N$ with an existential prime-power witness.
  \item The paper's probability distributions are described informally
    (``uniform over $\F_q^m$'', etc.); the \Lean{} formalization defines a
    custom \texttt{Distribution} type with explicit finite support.
  \item The \Lean{} formalization may include default instances for
    families whose outcomes are all zero, which the paper's prose never
    justifies.\footnote{See \texttt{MIPStarRE/LDT/Commutativity/Transport/FullSlice.lean},
      line~25, for the \texttt{zeroFullSliceOpFamily} definition.}
\end{enumerate}

The paper requires $q$ to be a prime power (Section~\ref{sec:prelims} of
\cite{Ji2020LowIndividualDegree}) because it uses \emph{finite fields}
throughout: the finite field trace $\tr : \F_q \to \F_p$, the additive
character $\omega^{\tr[\cdot]}$, and the Schwartz--Zippel lemma which
operates over the field $\F_q$. The field structure is needed for
polynomial evaluation, line questions, and the Fourier orthogonality facts
(\Cref{prop:fourier-fact-scalar,prop:fourier-fact-vector}).

\section{The point at issue}

\subsection{The field parameter $q$}

In the paper, $q$ is defined implicitly by ``$q = p^t$ is a prime power.''
The Lean \texttt{Parameters} structure\footnote{\texttt{MIPStarRE/LDT/Basic/ParametersBase.lean}, line~36.}
stores $q : \N$ with a field \texttt{hqPrimePower : $\exists$ p n, Nat.Prime p $\land$ 0 < n $\land$ q = p \^{} n}.
This is an existential without computational content: the actual field
operations are provided separately by the typeclass
\texttt{FieldModel params.q},\footnote{\texttt{ParametersBase.lean}, line~224.}
which bundles a field $K$ with $\Fintype$ structure and a coding
equivalence $K \simeq \Fin q$. A canonical instance is built from the
prime-power witness via \texttt{PrimePowerFieldSpec.toFieldModel},
producing an honest \texttt{GaloisField p n}.\footnote{Honest finite fields
  from Mathlib: \texttt{GaloisField} provides the unique field of order
  $p^n$.}

Hence the \Lean{} formalization separates the \emph{assertion} that $q$ is
a prime power (stored as a proof) from the \emph{use} of a finite field of
order $q$ (accessed through the \texttt{FieldModel} typeclass). The
algebraic operations on scalars are then provided by the field structure,
and the coding bridge \texttt{decodeScalar}/\texttt{encodeScalar}
transports them to the combinatorial carrier $\Fin q$. The prime-power
requirement is used to obtain $2 \le q$, nonzero cardinality, and the
honest field model, all of which are explicit in the paper.

This is a formalization design choice, not a mathematical gap. The paper's
field $\F_q$ is faithfully represented by the \texttt{FieldModel} typeclass.

\subsection{The Distribution type}

The paper uses probability distributions throughout: uniformity over
$\F_q^m$, uniformity over lines, and explicit distributions for sampled
questions. The Lean code defines a custom \texttt{Distribution}
type\footnote{\texttt{MIPStarRE/LDT/Basic/Distribution.lean}, line~17.}
carrying a finite support \texttt{Finset $\alpha$} and non-negative
real-valued weights, together with an \texttt{IsProbability} predicate
requiring total mass $1$. The uniform distribution on a nonempty finite
type is a derived case. The averaging combinator \texttt{avgOver} sums
over the explicit support.

This is a project-local implementation of a formal concept that Mathlib's
probability theory library could also provide. The \texttt{Distribution}
type is faithful to the paper's uniform distributions and to the explicit
finite averages that appear in the paper's estimates. The design does not
change any mathematical statement.

\subsection{Zero families}

The Lean code defines \texttt{zeroFullSliceOpFamily}\footnote{\texttt{MIPStarRE/LDT/Commutativity/Transport/FullSlice.lean}, line~26.}
as an \texttt{OpFamily} with all outcomes and total equal to
$0$. This family serves as a neutral element in the state-dependent
distance comparisons used in the full-slice commutation transport.
Its use is proved safe by the lemma
\texttt{fullSliceProductLeft\_to\_zero\_le\_one},\footnote{FullSlice.lean,
  line~216.} which bounds the SDD distance between the zero family and the
actual full-slice product by $1$ (on a normalized state).

The paper never defines a zero family; it instead writes explicit
inequalities from which the same bounds follow. The Lean zero family is a
convenience that packages those explicit bounds into the
\texttt{SDDOpRel} relation. There is no mathematical assertion in the
paper that the Lean zero family would contradict.

\section{Conclusion}

All three design choices are faithful to the cited source. The
field-model typeclass together with the prime-power witness captures
the paper's finite field $\F_q$ without changing any mathematical
hypothesis. The \texttt{Distribution} type provides a concrete
implementation of the paper's probability distributions that is
sufficient for all finite averages that appear. The zero family is an
internal formalization convenience whose safety is proved.

\noindent\textbf{Verdict:} no discrepancy. The three design choices are
formalization decisions that faithfully capture the paper's mathematics.

\bibliographystyle{alpha}
\bibliography{docs/paper-gaps/references}

\end{document}